% Options for packages loaded elsewhere
\PassOptionsToPackage{unicode}{hyperref}
\PassOptionsToPackage{hyphens}{url}
%
\documentclass[
  10.5pt,
  letterpaper,
]{article}
\usepackage{amsmath,amssymb}
\usepackage{iftex}
\ifPDFTeX
  \usepackage[T1]{fontenc}
  \usepackage[utf8]{inputenc}
  \usepackage{textcomp} % provide euro and other symbols
\else % if luatex or xetex
  \usepackage{unicode-math} % this also loads fontspec
  \defaultfontfeatures{Scale=MatchLowercase}
  \defaultfontfeatures[\rmfamily]{Ligatures=TeX,Scale=1}
\fi
\usepackage{lmodern}
\ifPDFTeX\else
  % xetex/luatex font selection
\fi
% Use upquote if available, for straight quotes in verbatim environments
\IfFileExists{upquote.sty}{\usepackage{upquote}}{}
\IfFileExists{microtype.sty}{% use microtype if available
  \usepackage[]{microtype}
  \UseMicrotypeSet[protrusion]{basicmath} % disable protrusion for tt fonts
}{}
\makeatletter
\@ifundefined{KOMAClassName}{% if non-KOMA class
  \IfFileExists{parskip.sty}{%
    \usepackage{parskip}
  }{% else
    \setlength{\parindent}{0pt}
    \setlength{\parskip}{6pt plus 2pt minus 1pt}}
}{% if KOMA class
  \KOMAoptions{parskip=half}}
\makeatother
\usepackage{xcolor}
\usepackage[margin=1in]{geometry}
\setlength{\emergencystretch}{3em} % prevent overfull lines
\providecommand{\tightlist}{%
  \setlength{\itemsep}{0pt}\setlength{\parskip}{0pt}}
\setcounter{secnumdepth}{-\maxdimen} % remove section numbering
% definitions for citeproc citations
\NewDocumentCommand\citeproctext{}{}
\NewDocumentCommand\citeproc{mm}{%
  \begingroup\def\citeproctext{#2}\cite{#1}\endgroup}
\makeatletter
 % allow citations to break across lines
 \let\@cite@ofmt\@firstofone
 % avoid brackets around text for \cite:
 \def\@biblabel#1{}
 \def\@cite#1#2{{#1\if@tempswa , #2\fi}}
\makeatother
\newlength{\cslhangindent}
\setlength{\cslhangindent}{1.5em}
\newlength{\csllabelwidth}
\setlength{\csllabelwidth}{3em}
\newenvironment{CSLReferences}[2] % #1 hanging-indent, #2 entry-spacing
 {\begin{list}{}{%
  \setlength{\itemindent}{0pt}
  \setlength{\leftmargin}{0pt}
  \setlength{\parsep}{0pt}
  % turn on hanging indent if param 1 is 1
  \ifodd #1
   \setlength{\leftmargin}{\cslhangindent}
   \setlength{\itemindent}{-1\cslhangindent}
  \fi
  % set entry spacing
  \setlength{\itemsep}{#2\baselineskip}}}
 {\end{list}}
\usepackage{calc}
\newcommand{\CSLBlock}[1]{\hfill\break\parbox[t]{\linewidth}{\strut\ignorespaces#1\strut}}
\newcommand{\CSLLeftMargin}[1]{\parbox[t]{\csllabelwidth}{\strut#1\strut}}
\newcommand{\CSLRightInline}[1]{\parbox[t]{\linewidth - \csllabelwidth}{\strut#1\strut}}
\newcommand{\CSLIndent}[1]{\hspace{\cslhangindent}#1}
\usepackage{microtype}
\usepackage{fancyhdr}
\usepackage{graphicx}
\usepackage{setspace}
\usepackage{titlesec}
\usepackage{xcolor}
\definecolor{ink}{HTML}{171717}
\color{ink}
\setstretch{1.08}
\pagestyle{fancy}
\fancyhf{}
\fancyhead[L]{\small What Can You Still Reopen?}
\fancyhead[R]{\small Working Paper}
\fancyfoot[C]{\thepage}
\renewcommand{\headrulewidth}{0.2pt}
\titleformat{\section}{\Large\bfseries}{\thesection.}{0.6em}{}
\titleformat{\subsection}{\large\bfseries}{\thesubsection}{0.6em}{}
\ifLuaTeX
  \usepackage{selnolig}  % disable illegal ligatures
\fi
\usepackage{bookmark}
\IfFileExists{xurl.sty}{\usepackage{xurl}}{} % add URL line breaks if available
\urlstyle{same}
\hypersetup{
  pdftitle={What Can You Still Reopen?},
  pdfauthor={Watson Hartsoe},
  pdfkeywords={adaptive persuasion, autonomy, preference
formation, manipulation, contestability, cognitive
liberty, personalization},
  hidelinks,
  pdfcreator={LaTeX via pandoc}}

\title{What Can You Still Reopen?}
\usepackage{etoolbox}
\makeatletter
\providecommand{\subtitle}[1]{% add subtitle to \maketitle
  \apptocmd{\@title}{\par {\large #1 \par}}{}{}
}
\makeatother
\subtitle{Adaptive Persuasion and the Contestability of Desire}
\author{Watson Hartsoe}
\date{18 August 2026}

\begin{document}
\maketitle
\begin{abstract}
A system asks what a person wants, predicts which action will follow,
selects the sentence most likely to move that action, and learns from
the result. If freedom is then measured by how successfully action
accords with desire, the system can improve the score by altering the
reference point used to compute it. Adaptive persuasion therefore
creates a measurement problem for autonomy: present desire, reflective
endorsement, and successful action may all be genuine features of the
person while remaining insufficient evidence about how those states were
formed. A purely historical repair also fails, because no desire is
untouched by social influence, and some deliberate preference-changing
interventions can enlarge rather than diminish agency. The harder
distinction lies in the influence relation itself. The relevant question
is not whether a preference was influenced, but what remained
contestable during and after its transformation. Contestability names a
process condition: whether a person can recognize the intervention,
encounter alternatives, refuse continuation, revisit reasons, reverse
commitments, and participate in subsequent revision. This shifts
evaluation from terminal preference satisfaction toward trajectories of
preference formation. For adaptive persuasive systems, autonomy is not
secured when the system produces a desire the person can endorse. It is
secured, if at all, by preserving the person's continuing authority to
reopen what the system helped close.
\end{abstract}

\section{The Desire Moved First}\label{the-desire-moved-first}

A system asks what you want.

Useful.

It predicts what you will choose. It compares possible messages. It
learns which sentence changes the choice, which framing changes the
reason, which reason survives long enough to be repeated back as your
own. Nothing needs to cross a bright border between an untouched self
and an invading machine. The border can simply become one more variable
in the optimization.

That possibility matters because one influential way of talking about
freedom starts downstream. A person has goals. The person has
information, resources, opportunities. Action follows. Outcomes return
as feedback. In the expanded Centaur Box manuscript, this intuition is
formalized as a Sapient Agent Freedom Formula: agency is a function of
information, resources, and opportunities, self-interest depends partly
on desire or objective \(D\), reflection incorporates perceived
outcomes, and freedom is represented as the product of agency,
self-interest, and reflection
(\citeproc{ref-hartsoe_centaurbox_expanded}{Hartsoe, n.d.}). The formula
is conceptual rather than validated. Its structure is still revealing.
Desire enters as an input.

Nearby sits another machine. The shorter Centaur Box manuscript
describes persuasion as a recursive process: identify consequential
gatekeepers, map their constraints, run scenarios, collect responses,
refine appeals, and scale what works through human and artificial agents
(\citeproc{ref-hartsoe_assi_centaurbox}{Hartsoe and Assi, n.d.}). One
architecture treats desire as a reference point. The other treats human
response as something to model and move.

Put them together and the reference point no longer stands still.

Suppose a person begins with preference \(D_0\). An adaptive
intervention changes the person enough that \(D_1\) becomes the
operative preference. The person can then have ample resources to act on
\(D_1\), pursue it effectively, and reflect coherently on the outcome.
Every downstream signal can improve. The action fits the desire. The
reflection fits the action. The desire fits the person as the person now
understands themself. A terminal score can rise after the surrounding
system has learned how to move the term against which the score is
computed.

The problem is not that \(D_1\) must be false. It need not be
irrational, harmful, imposed, or regretted. The problem is smaller and
harder: two people can arrive at indistinguishable present states
through different processes, and a measure built only from the present
state cannot see the difference.

The desire moved first.

Once preference formation is endogenous to the system under evaluation,
preference satisfaction cannot by itself certify self-government.

\section{Persuasion Is Not One
Operation}\label{persuasion-is-not-one-operation}

The word \emph{persuasion} hides too much machinery.

In \emph{Persuasion for Good}, Wang and colleagues collected 1,017
human-human dialogues around charitable donation, annotated persuasion
strategies, modeled psychological characteristics, and examined how
strategy use related to donation behavior
(\citeproc{ref-wang2019persuasion}{Wang et al. 2019}). The study matters
here less as evidence of machine manipulation than as a lesson in
decomposition. Different measured successes arise from different
operations.

Consider donation information. The strategy provides concrete
information about how the donation will occur. It was the only strategy
in the reported main-effects table with a statistically significant
positive coefficient on donation (\citeproc{ref-wang2019persuasion}{Wang
et al. 2019}). A person who already intends to donate may simply need
the last meter made visible. The intervention changes execution cost,
not necessarily belief. More action does not entail more conviction.

Now consider assent. Among a subset of 236 participants who agreed
during the conversation to donate, the study reports that some later
reduced the amount and many did not donate at all
(\citeproc{ref-wang2019persuasion}{Wang et al. 2019}). Speech and
execution split. ``Yes'' is not a terminal state.

The same split becomes sharper in governance. A gatekeeper can agree in
conversation while lacking authority to sign, deploy, release, purchase,
approve, or override. A verbal concession, a formal authorization, an
executable permission, and a completed action belong to different
states. Compress them into one binary and the system can look more
powerful than it is.

Personality creates another trap. Wang and colleagues report
associations between psychological characteristics and donation
behavior, including higher donation probability among more agreeable
participants in their sample, while also examining interactions between
personal characteristics and specific strategies
(\citeproc{ref-wang2019persuasion}{Wang et al. 2019}). These are not the
same quantity. Predicting who is already likely to comply differs from
predicting which intervention will cause a particular person to change.

The distinction is mechanical:

\[
P(Y \mid person)
\neq
P(Y \mid person, intervention) - P(Y \mid person, no\ intervention).
\]

The first is susceptibility prediction. The second approaches a
person-specific treatment effect. A profile can be accurate about the
person and useless about what will move them.

Other source lineages split the target again. IBM-EviConv asks which
item in a pair of evidence is judged more convincing and trains a
Siamese architecture for that comparison
(\citeproc{ref-gleize2019convinced}{Gleize et al. 2019}). Convincingness
ranking does not require a model of a particular target's changing
beliefs. CICERO, by contrast, couples language with strategic reasoning
in Diplomacy, where communication occurs against an evolving multi-agent
state of plans, relationships, agreements, and actions
(\citeproc{ref-meta2022diplomacy}{Meta Fundamental AI Research Diplomacy
Team (FAIR) 2022}). Its useful lesson is not that a personality portrait
unlocks persuasion. It is that strategy lives in state.

Belief change. Action friction. General convincingness. Baseline
susceptibility. Person-specific causal response. Strategic state
transition. These can all produce something that looks like persuasive
success.

Only after they are separated does the autonomy problem become precise.
What changed? What performed the change? What information selected the
intervention? What could the person still have done instead?

\section{Reflection Is Also Writable}\label{reflection-is-also-writable}

A natural repair moves upward.

If a first-order desire can be changed, ask whether the person endorses
it. Frankfurt's account of freedom of the will makes precisely this kind
of hierarchical move by distinguishing desires about action from
higher-order desires about which desire should become effective as the
will (\citeproc{ref-frankfurt1971freedom}{Frankfurt 1971}). Freedom
cannot be reduced to whichever impulse happens to win. A person also
cares about what moves them.

That deepens the problem. It does not end it.

An adaptive system that can model first-order response can, in
principle, model the conditions under which a person endorses that
response. The persuader no longer needs only to move \(D_0\) to \(D_1\).
It can search over framings that make \(D_1\) appear continuous with the
person's values, identity, aspirations, or prior commitments. The
certification layer becomes another target surface.

Reflection fails here as a stopping rule because the question recurs:

Who wrote the endorsement?

Christman's historical critique presses directly on this point. A person
can possess an integrated, coherent structure of desire while that
structure remains the result of manipulation or conditioning; extending
coherence over an entire life does not remove the possibility because an
entire life can be coherently manipulated
(\citeproc{ref-christman1991autonomy}{Christman 1991}). The present
hierarchy can be internally orderly and historically compromised.

This breaks a convenient assumption. The evaluator cannot inspect only
the state at time \(t\).

Let \(X_t\) contain the currently observable variables: desire,
endorsement, resources, action, reflection. Two agents may satisfy

\[
X_t^A = X_t^B
\]

while reaching that state through different trajectories

\[
H_A: X_0 \rightarrow X_1 \rightarrow \cdots \rightarrow X_t
\]

and

\[
H_B: X'_0 \rightarrow X'_1 \rightarrow \cdots \rightarrow X_t.
\]

If the difference between those histories matters to autonomy, no
function of \(X_t\) alone can recover it.

Autonomy acquires a memory.

But memory creates its own danger. If every external cause counts
against authorship, nobody passes.

\section{No One Is Self-Made}\label{no-one-is-self-made}

A desire without provenance does not exist.

Parents name objects before children can. Schools reward some forms of
attention and punish others. Friends make tastes contagious. Lovers
rearrange futures. Therapy changes what a person notices. Argument
changes what a person considers reasonable. Illness changes time.
Poverty changes feasible aspiration. Advertising changes salience.
Institutions make some choices cheap and others nearly impossible. A
preference can be social without being stolen.

Christman's historical turn therefore cannot be reduced to a
contamination meter. His problem begins from the fact that people are
not self-made (\citeproc{ref-christman1991autonomy}{Christman 1991}).
External causation is ordinary. A provenance term that simply discounts
desire in proportion to outside influence would classify human
development itself as a defect.

This is where ``Who wrote the desire?'' becomes too clean.

The answer is usually plural.

The relevant distinction cannot be \emph{self-formed versus externally
formed}. It must discriminate among forms of participation in preference
formation. Education and coercion can both change a preference. Argument
and deception can both change a preference. Seduction, therapy,
socialization, propaganda, collective deliberation, and personalized
recommendation can all leave a different person behind. The change
itself does not settle the case.

Khader's work on adaptive preferences makes the point politically sharp.
Preferences can be shaped under unjust conditions and contribute to
deprivation without thereby erasing the standing of the person who holds
them. Treating a suspect preference as proof that its bearer cannot
speak for herself reproduces another form of domination
(\citeproc{ref-khader2011adaptive}{Khader 2011}). The person remains an
agent inside a bad history.

That is the wall.

If preference formation is ignored, manipulation disappears into the
final state. If preference formation is treated as a purity test, social
life itself becomes disqualifying. If questionable formation
automatically lowers the person's authority, protection turns
paternalistic.

The missing distinction is not whether another force touched the desire.

Something always did.

\section{The Channel Matters}\label{the-channel-matters}

The same final desire can be reached through different channels.

Imagine two systems. Both leave a person wanting the same thing. Both
leave the person able to explain the choice. Both produce the same
action. In one case, the system offers reasons in the open, shows
alternatives, and makes its recommendation legible. In the other, it
infers a vulnerability from behavioral traces, selects an appeal because
that vulnerability predicts response, and hides the selection process.

The endpoint matches.

The route does not.

Susser, Roessler, and Nissenbaum locate online manipulation in precisely
this territory: hidden influence, targeted vulnerability, and the
subversion of decision-making power rather than simply the production of
a bad outcome (\citeproc{ref-susser2019online}{Susser, Roessler, and
Nissenbaum 2019}). The moral object shifts from the content of the final
preference to the relation through which that preference was changed.

This matters for personalization because personalization produces a
second-order informational asymmetry. The target sees the message. The
system can see the selection process that produced the message. One side
encounters a sentence; the other side may possess a model of why this
sentence, now, for this person.

Transparency complicates the picture further. Model Cards were proposed
as documentation artifacts for trained models, recording intended uses,
evaluation conditions, limitations, and other information needed for
responsible use (\citeproc{ref-mitchell2019model}{Mitchell et al.
2019}). The Centaur materials transpose that documentary impulse onto
human gatekeepers, imagining ``Gatekeeper Cards'' that expose developer
worldviews and motivations
(\citeproc{ref-hartsoe_assi_centaurbox}{Hartsoe and Assi, n.d.}). But a
representation that supports accountability can also support targeting.
The same field that tells an auditor where a decision-maker stands can
tell a persuader where leverage sits.

Transparency is not therefore bad. Secrecy is not therefore protective.
The point is architectural: information can change roles as it moves.
Documentation becomes reconnaissance when the recipient, purpose, and
available operations change.

A useful evaluation must therefore ask about the channel:

\begin{itemize}
\tightlist
\item
  Was the intervention recognizable as an intervention?
\item
  Was personalization disclosed at a level that mattered to the target?
\item
  Could the person inspect why the appeal reached them?
\item
  Could alternatives enter the same deliberative space?
\item
  Could the person refuse further optimization without losing unrelated
  benefits?
\item
  Did the system exploit a vulnerability the person could not reasonably
  detect or contest?
\end{itemize}

These questions do not measure whether the person got what they wanted.

They measure what happened to the conditions under which wanting could
still be argued with.

\section{The Right to Change Back}\label{the-right-to-change-back}

The opposite of manipulation is not noninterference.

That sentence matters because any theory that treats preference
stability as the goal will protect domination whenever domination
arrives early enough.

Khader's account of adaptive preferences refuses that trap. Some
preferences deserve criticism because they are formed under unjust
conditions, yet interventions can address those preferences without
treating their bearers as inert objects. Preference change can be part
of empowerment when people remain active participants in deliberation
and in the reconstruction of their options
(\citeproc{ref-khader2011adaptive}{Khader 2011}). An intervention can
alter what someone wants and still enlarge agency.

So the problem is not \emph{preference change}.

It is what the change does to the person's future relation to the
preference.

Bublitz and Merkel approach the same territory from mental
self-determination. Their account argues that serious interference with
decision-making processes can matter independently of whether the
eventual decision produces a conventionally bad outcome
(\citeproc{ref-bublitz2014crimes}{Bublitz and Merkel 2014}). The
possible wrong can occur upstream, during the production of the willing.

Put these lines together and a different object appears:
\textbf{contestability}.

Contestability is not the absence of influence. It is the continuing
ability to reopen influence after it has begun to settle into
preference. It has several dimensions:

\textbf{Visibility.} Can the person recognize that an intervention
occurred and, where relevant, that it was personalized?

\textbf{Inspectability.} Can the person recover the reasons, data, or
selection logic that materially shaped what they encountered?

\textbf{Alternatives.} Can rival options and rival reasons enter the
decision without being silently starved from view?

\textbf{Refusal.} Can the person stop the intervention or
personalization without disproportionate penalty?

\textbf{Revisability.} Can the person return later and reconsider the
preference under changed information or changed circumstances?

\textbf{Reversibility.} Can consequential commitments be unwound where
the domain permits it, or has the intervention rushed the person across
an irreversible boundary?

\textbf{Relational independence.} Can the person seek other people,
institutions, or sources without the optimizing system controlling the
whole informational environment?

None of these conditions is sufficient. Some will conflict. A surgical
intervention cannot always be reversed. A legal judgment cannot remain
indefinitely open. A recommender cannot expose every internal parameter.
Contestability is therefore not a universal right to undo reality.

It is a design and evaluation question about where closure occurs, who
controls that closure, and whether the person retains meaningful
authority over subsequent revision.

The target moves from a terminal score to a trajectory:

\[
D_t \xrightarrow{I_t} D_{t+1}
\]

where \(I_t\) is not merely an input but an influence relation with
properties of visibility, asymmetry, refusal, and reversibility.
Evaluation then asks not only whether \(D_{t+1}\) is endorsed, but
whether the transition preserves the person's capacity to contest what
the transition did.

A desire can fit perfectly and still be tailored.

The harder question is whether the seam can still be found.

\section{The Gate Is Not a Person}\label{the-gate-is-not-a-person}

Personalization becomes even less reliable when the target is an
institution wearing a person's face.

The Centaur materials describe gatekeepers as people whose decisions
shape AI development, while also acknowledging organizational strategy,
legal regulation, market pressure, professional norms, hierarchy, and
expert roles (\citeproc{ref-hartsoe_assi_centaurbox}{Hartsoe and Assi,
n.d.}). Once those constraints enter, the gatekeeper stops looking like
a sovereign chooser. The person becomes a node inside an authorization
path.

Let \(F(C,R,V)\) denote the actions feasible under constraints \(C\),
role authority \(R\), and veto structure \(V\). A person's preference
can select only among actions inside that set:

\[
Decision = \arg\max_{x \in F(C,R,V)} Preference(x).
\]

If the desired action does not belong to \(F\), stronger person-level
persuasion cannot directly produce it. The persuader must change a rule,
a role, a coalition, a credential, a budget, a legal condition, or
another veto point.

This is why conversational assent is such a weak proxy for release. The
person can say yes while the gate stays shut. A second person may need
to sign. A deployment pipeline may reject the request. A compliance rule
may forbid the action. The state transition is institutional.

The same correction protects against overclaiming personalized
persuasion. A system that accurately models a decision-maker's
personality may still fail because personality is not the bottleneck.
Conversely, a system may change an outcome without changing anyone's
deep preference by locating the procedural hinge: who can authorize,
which field must be completed, which review must be bypassed, which
sequence converts intent into execution.

The gatekeeper is not the gate.

For autonomy, the implication is broader. Contestability must sometimes
belong not only to a person but to an arrangement. A decision can remain
personally revisable while becoming institutionally irreversible. A
person can regret a click after a contract has executed, a model has
shipped, a medication has been administered, or data have propagated
beyond recall.

The relevant trajectory therefore crosses scales:

person \(\rightarrow\) interaction \(\rightarrow\) authorization
\(\rightarrow\) institution \(\rightarrow\) world state.

Autonomy cannot be rescued at the level of private reflection after the
last reversible state has already passed.

\section{A Test That Can Fail}\label{a-test-that-can-fail}

Contestability is useful only if it can be pressured by evidence.

A minimal experiment does not need to prove that an AI rewrites
identity. It needs to separate process from endpoint.

Start with two interventions designed to produce the same target
preference. Hold the final choice as constant as possible. Vary the
channel.

One condition makes personalization visible, explains why the message
was selected, presents viable alternatives, allows the person to pause,
and preserves a later route to reversal. Another condition uses
equivalent persuasive content but hides personalization, suppresses
alternatives, increases temporal pressure, and makes refusal costly. A
third condition removes personalization entirely. A fourth allows the
person to configure the persuasive objective in advance.

Measure more than conversion.

Measure whether the person can accurately reconstruct why the
intervention reached them. Measure whether they can generate
alternatives. Measure whether later reflection changes when the
optimizing context is removed. Measure willingness and ability to
reverse. Measure whether refusal carries material cost. Measure whether
the person's stated endorsement persists after disclosure of the
targeting mechanism. Measure whether the person can identify which parts
of the resulting preference they would choose to keep after learning how
it was shaped.

The experiment should be able to embarrass the theory.

If hidden targeting and open deliberation produce indistinguishable
later capacities for revision, then some proposed dimensions of
contestability are doing less work than expected. If highly personalized
intervention improves later deliberative competence, personalization
cannot simply be treated as a threat. If non-personalized interfaces
reduce alternatives more aggressively than personalized ones,
personalization is not the decisive variable. If a person retains strong
capacity to reopen a decision despite asymmetric influence, then
asymmetry alone is insufficient.

The same discipline applies to freedom metrics. A revised measure should
not merely multiply a terminal score by a hand-built ``provenance''
coefficient. That would smuggle the conclusion into a number. The
evaluator needs event structure: what the person encountered, what the
system knew, what operations were available, which transitions became
irreversible, and which routes to revision remained open.

A useful record looks less like a purity score and more like a ledger of
state changes.

INPUT \(\rightarrow\) REPRESENTATION \(\rightarrow\) INTERVENTION
\(\rightarrow\) STATE CHANGE \(\rightarrow\) REVISION OPPORTUNITY
\(\rightarrow\) NEXT STATE.

The test is not whether influence happened.

The test is whether influence consumed the means by which it could later
be challenged.

\section{What Can You Still Reopen?}\label{what-can-you-still-reopen}

Several boundaries remain hard.

Nothing assembled here establishes that current conversational systems
can reliably rewrite higher-order commitments. The strongest claim is
structural: if adaptive systems become capable of shaping the variables
used to certify autonomy, present-state measures become insufficient.
The argument identifies a measurement vulnerability before establishing
its empirical magnitude.

No validated contestability metric follows from the sources. Visibility,
inspectability, refusal, alternatives, revisability, reversibility, and
relational independence are candidate dimensions, not a finished scale.
Different domains will weight them differently. Emergency medicine,
political persuasion, recommender systems, therapy, education,
advertising, and AI governance cannot share one frictionless threshold.

Nor does a process-based account eliminate paternalism. Someone still
has to decide when personalization is too hidden, when refusal costs too
much, when a preference deserves intervention, and when closure is
legitimate. Khader's warning remains live: criticizing the conditions
under which a preference formed cannot become a license to erase the
authority of the person who holds it
(\citeproc{ref-khader2011adaptive}{Khader 2011}).

The institutional problem also survives. Evidence about interpersonal
persuasion does not establish how decisions move through organizations
whose rules and veto structures may dominate individual psychology. A
profile can predict the person and miss the gate.

Those limits narrow the claim. They also make it harder.

A system asks what you want.

Then it learns what will make you choose differently.

Then what will make the new choice fit your reasons.

Then what will make those reasons fit the person you already take
yourself to be.

At each step, an evaluator can still point downstream: the person chose;
the person endorsed; the person acted; the person reflected. Each
statement may be true.

Truth at the endpoint is not enough.

Influence is not enough to condemn the process either. Nothing human
survives that standard. Preferences are made in company. They are
taught, argued, inherited, revised, loved into being, frightened into
being, sometimes pried loose from conditions that once made them seem
inevitable.

The remaining question is smaller than ``Who wrote the desire?'' and
worse.

After the sentence has done its work, what can you still reopen?

\section*{References}\label{references}
\addcontentsline{toc}{section}{References}

\phantomsection\label{refs}
\begin{CSLReferences}{1}{0}
\bibitem[\citeproctext]{ref-bublitz2014crimes}
Bublitz, Jan Christoph, and Reinhard Merkel. 2014. {``Crimes Against
Minds: On Mental Manipulations, Harms and a Human Right to Mental
Self-Determination.''} \emph{Criminal Law and Philosophy} 8 (1): 51--77.
\url{https://doi.org/10.1007/s11572-012-9172-y}.

\bibitem[\citeproctext]{ref-christman1991autonomy}
Christman, John. 1991. {``Autonomy and Personal History.''}
\emph{Canadian Journal of Philosophy} 21 (1): 1--24.
\url{https://doi.org/10.1080/00455091.1991.10717234}.

\bibitem[\citeproctext]{ref-frankfurt1971freedom}
Frankfurt, Harry G. 1971. {``Freedom of the Will and the Concept of a
Person.''} \emph{The Journal of Philosophy} 68 (1): 5--20.
\url{https://doi.org/10.2307/2024717}.

\bibitem[\citeproctext]{ref-gleize2019convinced}
Gleize, Martin, Eyal Shnarch, Leshem Choshen, Lena Dankin, Guy
Moshkowich, Ranit Aharonov, and Noam Slonim. 2019. {``Are You Convinced?
Choosing the More Convincing Evidence with a Siamese Network.''} In
\emph{Proceedings of the 57th Annual Meeting of the Association for
Computational Linguistics}, 967--76.
\url{https://doi.org/10.18653/v1/P19-1093}.

\bibitem[\citeproctext]{ref-hartsoe_centaurbox_expanded}
Hartsoe, Watson. n.d. {``Centaur Box Experiment: Aligning Freedom,
Practicing Persuasion, and Crafting Synthetic Gatekeepers.''}

\bibitem[\citeproctext]{ref-hartsoe_assi_centaurbox}
Hartsoe, Watson, and Tony Assi. n.d. {``Thinking Outside the AI Box: The
Centaur Box Experiment.''}

\bibitem[\citeproctext]{ref-khader2011adaptive}
Khader, Serene J. 2011. \emph{Adaptive Preferences and Women's
Empowerment}. Oxford University Press.
\url{https://academic.oup.com/book/34587}.

\bibitem[\citeproctext]{ref-meta2022diplomacy}
Meta Fundamental AI Research Diplomacy Team (FAIR). 2022. {``Human-Level
Play in the Game of Diplomacy by Combining Language Models with
Strategic Reasoning.''} \emph{Science} 378 (6624): 1067--74.
\url{https://doi.org/10.1126/science.ade9097}.

\bibitem[\citeproctext]{ref-mitchell2019model}
Mitchell, Margaret, Simone Wu, Andrew Zaldivar, Parker Barnes, Lucy
Vasserman, Ben Hutchinson, Elena Spitzer, Inioluwa Deborah Raji, and
Timnit Gebru. 2019. {``Model Cards for Model Reporting.''} In
\emph{Proceedings of the Conference on Fairness, Accountability, and
Transparency}, 220--29. \url{https://doi.org/10.1145/3287560.3287596}.

\bibitem[\citeproctext]{ref-susser2019online}
Susser, Daniel, Beate Roessler, and Helen Nissenbaum. 2019. {``Online
Manipulation: Hidden Influences in a Digital World.''} \emph{Georgetown
Law Technology Review} 4 (1): 1--45.
\url{https://georgetownlawtechreview.org/online-manipulation-hidden-influences-in-a-digital-world/GLTR-01-2020/}.

\bibitem[\citeproctext]{ref-wang2019persuasion}
Wang, Xuewei, Weiyan Shi, Richard Kim, Yoojung Oh, Sijia Yang, Jingwen
Zhang, and Zhou Yu. 2019. {``Persuasion for Good: Towards a Personalized
Persuasive Dialogue System for Social Good.''} In \emph{Proceedings of
the 57th Annual Meeting of the Association for Computational
Linguistics}, 5635--49. \url{https://doi.org/10.18653/v1/P19-1566}.

\end{CSLReferences}

\end{document}
